\section{Parallel self-play in Rust}
\label{sec:parallel}

\subsection{What is actually scarce}
\label{sec:scarce}

A self-play system has two kinds of work, and they differ by two orders of
magnitude in cost. The tree work -- descending a path of ten edges, comparing a
few hundred weighted sums, playing and unplaying moves -- is measured in
microseconds. One evaluation of a small residual network is roughly $1.7$
gigarithmetic operations for the configuration of Table~\ref{tab:scaling}, which
on a device delivering a few effective tera-operations per second is measured in
milliseconds when it is done alone and much less when it is done in a batch.

Three consequences follow, and they determine the entire parallel design.

\emph{Throughput is set by evaluations per second.} Games per hour is
\begin{equation}
\text{games/h} \;=\; \frac{\text{evals/s}}{\text{simulations/move}}
\cdot \frac{1}{\text{moves/game}} \cdot 3600 ,
\label{eq:throughput}
\end{equation}
and neither the quality of the tree implementation nor the number of CPU cores
appears in it. Everything else in this section exists to raise the first factor.

\emph{A batch of one wastes the device.} Accelerator efficiency is a function of
batch size: the fixed cost of launching work, and the fact that the device is
built for wide parallel operations, mean that evaluating 128 positions in one
call costs far less than 128 calls of one. Batching is therefore not an
optimization; it is the difference between using the device and merely occupying
it.

\emph{But batching is latency, and latency starves the workers.} A worker that
submits a request and waits cannot do anything else. If the service waits for
more requests before running, every worker in the batch waits longer, and if the
wait exceeds the useful work per evaluation, the workers become the bottleneck
instead of the accelerator.

The design problem is therefore a single trade-off: \emph{maximize batch size
without making workers wait longer than the device needs}. The classical answer
in self-play is to build many independent trees, because more trees means more
outstanding requests, which means larger batches.

\subsection{The parallelization taxonomy, and why the choice is inverted}
\label{sec:taxonomy}

Parallel MCTS divides into three families \citep{chaslot2008parallel,
browne2012survey}, and it is worth stating why two of them do not apply here.

\paragraph{Leaf parallelization.} Run many simulations from one expanded leaf
concurrently. The technique exists to parallelize \emph{rollouts}: several random
playouts from the same leaf, averaged. An AlphaZero-style search has no rollouts
(Section~\ref{sec:algorithm}), so there is nothing to parallelize. Evaluations
are already batched by the service, which is the useful half of the idea.

\paragraph{Root parallelization.} Build several independent trees from the same
root and merge their root statistics at the end. Synchronization is zero. The
cost is that each tree repeats the exploration the others perform, so $k$ trees
do not search $k$ times as much of the tree; they search the same part $k$ times
and diversify only through randomness. With a shared batched evaluator this
family becomes attractive again, for a reason its original formulation did not
have: independent trees generate independent evaluation requests, which is
exactly what the batcher wants.

\paragraph{Tree parallelization.} One shared tree, many threads descending it.
The classical difficulty is that concurrent threads select the same best path and
converge on the same leaf, so a mechanism is needed to keep them apart: a
\emph{virtual loss}, in which a thread records a temporary loss on the edges it is
descending, deterring followers until the real value is backed up
\citep{chaslot2008parallel,segal2011paralleluct}. AlphaGo (2016) ran precisely
this construction, with lock-free tree updates and virtual loss
\citep{silver2016alphago}.

The historical reason is worth internalizing, because it explains why the same
choice is wrong here. \emph{AlphaGo needed tree parallelization because its
network evaluations were scarce and its search was deep.} A large network
evaluated slowly, and the search needed thousands of simulations per move, so the
only way to fill the device was to have many threads deepen one tree.
Self-play on a single accelerator is the opposite situation. Evaluations come
back in batches from one device; the number of simulations per move is a few
hundred; and what the batcher lacks is not threads inside one tree but
\emph{trees}.

Hence the choice of \emph{game-level parallelism}: each worker plays its own game
with its own single-threaded tree, its own board, and its own random-number
generator. This is root parallelization taken to the point where no merging
happens at all, which is correct here because the trees belong to different games
and produce different training data rather than different estimates of one move.
The properties that follow are the ones that matter:

\begin{itemize}
  \item \textbf{No shared tree, no locks.} A worker's data structures are private,
        so an entire class of concurrency defects cannot occur, and cache lines
        are not fought over.
  \item \textbf{The only synchronization point is the evaluator}, which is
        precisely where synchronization is wanted, since batching \emph{is} the
        coordination.
  \item \textbf{Scaling is linear in workers until the device saturates}, because
        the workers do not contend with each other at all.
\end{itemize}

The same choice was made by ELF OpenGo at datacenter scale: 32 self-play workers
per GPU, with eight game threads per worker and no shared trees
\citep{tian2019elf}. When a design that is forced by a single accelerator matches
what an industrial system chose with hundreds, the reasoning is likely sound.

\begin{figure}[t]
\centering
\begin{tikzpicture}[
  node distance=6mm and 8mm,
  w/.style={draw, rounded corners=2pt, align=center, font=\scriptsize,
            minimum width=26mm, minimum height=12mm, fill=black!3},
  svc/.style={draw, rounded corners=2pt, align=center, font=\small,
              minimum width=78mm, minimum height=13mm, fill=blue!4},
  buf/.style={draw, rounded corners=2pt, align=center, font=\small,
              minimum width=78mm, minimum height=10mm, fill=green!4},
  arr/.style={-{Latex[length=2mm]}, thick},
  darr/.style={{Latex[length=2mm]}-{Latex[length=2mm]}, thick},
]
\node[w] (w0) {worker 0\\game $+$ tree $+$ RNG};
\node[w, right=of w0] (w1) {worker 1\\game $+$ tree $+$ RNG};
\node[right=of w1, font=\large] (dots) {$\cdots$};
\node[w, right=of dots] (wN) {worker $W{-}1$\\game $+$ tree $+$ RNG};
\node[svc, below=10mm of w1, xshift=11mm] (svc)
  {\textbf{evaluator service} (one thread, owns the model)\\
   collect $\to$ batch ($\le B$ requests, $\le t$ window) $\to$ one forward pass $\to$ reply};
\node[buf, below=7mm of svc] (buf)
  {\textbf{replay window}: sharded ring buffer of games};
\draw[darr] (w0.south) -- ++(0,-4mm) -| ([xshift=-28mm]svc.north);
\draw[darr] (w1.south) -- (svc.north);
\draw[darr] (wN.south) -- ++(0,-4mm) -| ([xshift=28mm]svc.north);
\draw[arr] (svc.south) -- node[right,font=\scriptsize]{finished games} (buf.north);
\end{tikzpicture}
\caption{The actor layout. Workers never touch the model, the model never sees a
tree, and the service is the only place where threads interact.}
\label{fig:actors}
\end{figure}

\subsection{The evaluator service}
\label{sec:service}

The service is one thread that owns the model and runs a loop. In outline:

\begin{lstlisting}[style=rust,caption={The batching loop. One device call per
batch, one reply per request.},label={lst:service}]
loop {
    let first = rx.recv()?;                       // block until work exists
    let mut batch = vec![first];
    while batch.len() < MAX_BATCH {               // fill the batch
        match rx.recv_timeout(WINDOW) {           // ... within a time window
            Ok(req) => batch.push(req),
            Err(_)  => break,                     // window closed: run now
        }
    }
    let out = model.forward(&batch);              // ONE device call
    for (req, res) in batch.into_iter().zip(out) {
        let _ = req.reply.send(res);              // one reply per request
    }
}
\end{lstlisting}

Five decisions are embedded in that short loop.

\paragraph{Dynamic batching with a window.} The batch is bounded by both a count
and a time. The window exists to catch the tail of the arrival process, not to
build the batch: when many workers are blocked, the batch fills to the cap
immediately and the window contributes nothing to latency. This is the point to
get right when tuning, because the two failure modes are opposite. A window that
is too long adds latency to every request without improving the batch; a window
that is too short fails to gather the requests that are already available.

A simple model makes the trade-off visible. Suppose the service takes
$\tau(b) = \tau_0 + b\,\tau_1$ for a batch of $b$ (a fixed launch cost plus a
per-item cost), and suppose $Q$ workers are blocked waiting. If $Q \ge B$ the
batch is full, throughput is $B/\tau(B)$, and the window does not matter. If
$Q < B$ the service must decide between waiting up to $t$ to gather more requests
and running the smaller batch. Waiting is only worthwhile if the arrivals during
the wait more than compensate for the delay they add to the requests already
collected. In practice the useful rule is: \emph{size the batch cap so that a
healthy worker pool can fill it, and keep the window short}. The cap, not the
window, is the throughput knob.

\paragraph{One model, one thread.} The network lives exclusively on the service
thread. Workers send plain request structs and receive plain result structs; they
never hold a handle to a tensor, a device, or a module. This has a larger
consequence than it first appears.

In Rust terms, it means the system uses \emph{ownership} to make a class of
concurrency bugs unrepresentable, rather than \code{Arc<Mutex<Model>{}>} to manage
them at runtime. There is no reader--writer lock on the weights, so there is no
lock contention on the hottest resource; there is no question about whether
accelerator handles are \code{Send}, because none of them ever cross a thread
boundary; and replacing the model after a training phase is a single assignment
on a single thread, with no coordination protocol. The synchronization that
remains is exactly the batching that the design wants.

This is the most Rust-specific decision in the whole system, and it generalizes:
when two components must share a resource, first ask whether the resource can be
\emph{owned} by one of them and accessed by messages.

\paragraph{Replies travel inside requests.} Each request carries a one-shot
sender for its own reply. The service needs no table mapping request identifiers
to waiting parties, no correlation logic, and no cleanup for abandoned requests:
if a worker stops waiting, the send simply fails and the error is ignored. The
message is self-addressed, so the protocol cannot get out of step.

\paragraph{Bounded channels are backpressure.} The request channel is bounded. If
the device falls behind, workers block in \code{send} rather than accumulating an
unbounded queue of stale requests, and the system throttles itself to device
speed with no watchdog, no adaptive rate control, and no measurement. Note the
interaction with the batching loop: a bounded channel plus a
\code{recv\_timeout} window means the service is never starved by an idle pool
and never flooded by an eager one.

\paragraph{Errors need a policy.} If the service thread panics, every worker is
blocked in \code{recv} forever and the process hangs rather than failing. Two
cheap defences are worth building: the reply channel is dropped when the service
dies, so \code{recv} on the worker side returns an error that can be propagated;
and the run loop treats a dead service as a fatal condition with a message, not
as a retryable one.

\subsection{Training, inference, and the single-device problem}
\label{sec:phased}

One accelerator must serve two clients: the evaluator service during self-play,
and the trainer during the training phase. A system can either phase them or run
them concurrently.

\emph{Continuous} operation is what AlphaZero did, on separate hardware for
self-play and training. ELF OpenGo ran it continuously on shared hardware and
reported a large throughput advantage for the asynchronous pipeline, at the cost
of two specific complications: games whose moves were produced by different
network versions, which makes the training data heterogeneous, and batch
normalization statistics that lag behind the current weights, which they fixed by
recomputing those statistics from a sample of batches at intervals
\citep{tian2019elf,ioffe2015batchnorm}.

\emph{Phased} operation alternates: generate $N$ self-play games with the
evaluator alone; run $K$ training steps with the trainer alone; play evaluation
games; repeat. The reasoning for choosing this first is straightforward. On a
single device, concurrency is not really concurrency: two clients contend on one
queue, the driver serializes them anyway but less predictably than the scheduler
would, and the batch shapes fight each other, because training wants batches of
a thousand and inference wants batches of at most a few hundred. Phasing costs
nothing in correctness, makes every number reproducible per iteration, and keeps
the first working version debuggable, which matters more than a percentage of
throughput at the point where the system has never run.

The upgrade path stays open by construction, and this is the design argument for
phasing: once the evaluator and the trainer are separate owners of separate
resources, moving to continuous operation is a scheduling change, not a rewrite.
The decision rule to write down is: go continuous when profiling shows the phase
boundaries costing more than a fixed fraction of throughput, and pay ELF's
complications deliberately, including a minimum replay size before training
starts.

\subsection{The replay window}
\label{sec:buffer}

The buffer is the only mutable structure shared between threads, so it should be
as simple as possible, and it should be sized by an argument rather than by
intuition.

The argument is about what to store. Storing \emph{positions} means storing the
encoded planes, the target distribution, and the outcome: about a kilobyte per
sample, simple, and immediately usable. Storing \emph{games} means storing the
move list and the metadata: tens of bytes per game, because the planes are
derivable. The difference is a factor of several thousand -- two million games
occupy on the order of a hundred megabytes of move lists, against terabytes for
the equivalent positions.

Storing games is the better choice for three reasons. It fits comfortably in
memory without paging. It keeps the encoding decisions mutable: if the plane
layout changes during experimentation, every stored game remains valid, whereas
stored planes are frozen at the old layout. And it puts the encoding work where
the CPU is idle: the trainer thread is device-bound, not CPU-bound, so replaying
a game through the rules engine at sampling time is nearly free, and it is the
same discipline that applies to any data pipeline -- store the raw form, encode
at the boundary.

Concurrency for the buffer follows the sharding principle already established.
Appends arrive from many workers and sampling happens on one thread. A single
mutex, or a lock-free structure, is over-engineering: a sharded ring buffer with
one mutex per shard keeps lock hold times at the level of a few instructions,
the sampler round-robins the shards, and the worst case is a slightly uneven
distribution of games rather than a correctness problem. Contention here is
measured in microseconds per game, not per sample.

Persistence should be per iteration rather than per sample: appending each
finished game to a journal file bounds a crash to the current phase, at one write
per game. A system that trains for days and loses the run to a power failure will
be rebuilt with a journal after the first incident, which is one incident later
than necessary.

\subsection{Sizing, and the one number to watch}
\label{sec:sizing}

Worker count follows from cores, not from a formula. Reserve one core for the
service thread, which is mostly idle waiting on the device but must never be
scheduled late, because a late service thread stalls every worker; reserve one
for the trainer and the operating system; give the rest to workers. On an
18-core machine that is 14 workers, and it is a starting value.

The tuning signal is a single number: \emph{evaluator queue depth}, the number of
requests waiting when the service starts a batch.

\begin{itemize}
  \item Queue consistently starving (batch sizes far below the cap): add workers,
        or lower the batch cap so that the batches that do exist complete sooner.
  \item Queue consistently flooded (the bounded channel is full and workers are
        blocking in \code{send}): the device is saturated, and extra workers only
        add latency. Remove workers, or reduce simulations per move.
\end{itemize}

This is the entire control loop, and it is deliberately a measurement rather than
a model. A worked estimate of the target throughput is useful for planning and
should be labelled as derived rather than measured. With the network of
Table~\ref{tab:scaling} at roughly $1.7$ gigarithmetic operations per evaluation,
a conservative assumption of five effective tera-operations per second gives
about $2900$ evaluations per second. At $400$ simulations per move that is about
seven moves per second across all workers, so a 50-move game completes in about
seven seconds and the system produces a few hundred games per hour. A
25{,}000-game iteration then takes on the order of two days. The precise numbers
depend on the device, and the useful conclusion is the order of magnitude: plan
for days to a first non-embarrassing agent and weeks to a strong one, and
distrust any plan more precise than that before the first measurement exists.

\subsection{Determinism}
\label{sec:determinism}

Self-play is stochastic in four places: the noise vector at the root, the
temperature sampling, the symmetry transform applied to training positions, and
the mini-batch sampling. Each should be driven by a generator that is owned by
exactly one worker, seeded from a run-level seed plus the worker index, so that
the $\text{worker} \times \text{game}$ stream can be reproduced exactly.

What can be promised depends on the backend. On a CPU backend, arithmetic is
deterministic and a fixed seed reproduces a run bit for bit. On an accelerator
backend, kernel selection and autotuning can make even the same operation produce
slightly different rounding between processes, so the promise is best-effort and
must be stated as such. This matters for one practical reason: reproducing a
strange game is by far the fastest route to diagnosing a search defect, and
reproducibility has to be designed in before it is needed. Recording the seed,
the mode, and the backend in the run journal costs nothing and converts an
unreproducible anomaly into an investigation.
